\subsection{RC Low Pass Filter}

\subsubsection{Overview}
Constructed from a resistor and capacitor in series, a RC
Low Pass Filter allows low frequency signals to pass while
attenuating high frequency signals. It is a single pole
filter but can be used in multiple stages to create higher
order filters.
\subsubsection{Schematic}
\begin{center}
  \begin{circuitikz}
    % Paths, nodes and wires:
	  \draw (5.5, 5.25) to[american voltage source, l_={$V_{ee}$}] (5.5, 6.5);
	  \node[ground] at (5.5, 8){};
	  \node[ground, xscale=-1, yscale=-1] at (5.5, 6.5){};
	  \draw (7, 9.25) to[european resistor, l={$G_1$}] (7, 7.75);
	  \draw (7, 6.75) to[capacitor, l={$C_1$}] (7, 5.25);
	  \draw (7, 7.75) -- (7, 6.75);
	  \draw (7, 5.25) -| (7, 5);
	  \draw (5.5, 5) -- (7, 5);
	  \draw (5.5, 5) -| (5.5, 5.25);
	  \draw (5.5, 9.25) -| (5.5, 9.5);
	  \draw (5.5, 9.5) -- (7, 9.5);
	  \draw (7, 9.5) -| (7, 9.25);
	  \draw (7, 7.25) -- (7.75, 7.25);
	  \draw (5.5, 8) to[sinusoidal voltage source, l_={$V_{in}$}] (5.5, 9.25);
	  \node[shape=rectangle, minimum width=0.965cm, minimum height=0.465cm](N1) at (8.25, 7.25){} node[anchor=center] at (N1.text){$V_{out}$};
	  \node[ground] at (9.5, 8){};
	  \node[ground, xscale=-1, yscale=-1] at (9.5, 5.25){};
	  \draw (11, 9.25) to[european resistor, l={$G_1$}] (11, 7.75);
	  \draw (11, 6.75) to[capacitor, l={$C_1$}] (11, 5.25);
	  \draw (11, 7.75) -- (11, 6.75);
	  \draw (11, 5.25) -| (11, 5);
	  \draw (9.5, 5) -- (11, 5);
	  \draw (9.5, 5) -| (9.5, 5.25);
	  \draw (9.5, 9.25) -| (9.5, 9.5);
	  \draw (9.5, 9.5) -- (11, 9.5);
	  \draw (11, 9.5) -| (11, 9.25);
	  \draw (11, 7.25) -- (11.75, 7.25);
	  \draw (9.5, 8) to[sinusoidal voltage source, l_={$V_{in}$}] (9.5, 9.25);
	  \node[shape=rectangle, minimum width=0.965cm, minimum height=0.465cm](N2) at (12.25, 7.25){} node[anchor=center] at (N2.text){$V_{out}$};
	  \node[shape=rectangle, minimum width=3.215cm, minimum height=0.465cm] at (6.125, 4.25){} node[anchor=north, align=center, text width=2.827cm, inner sep=6pt] at (6.125, 4.5){Full Schematic};
	  \node[shape=rectangle, minimum width=2.715cm, minimum height=0.465cm] at (10.375, 4.25){} node[anchor=north, align=center, text width=2.327cm, inner sep=6pt] at (10.375, 4.5){Small-Signal\\Schematic};
  \end{circuitikz}
\end{center}

\subsubsection{Transfer Function}
Applying KCL at $V_{out}$ on Full Schematic, we get:
\[
  (V_{out} - V_{in}){G_1} + (V_{out} - V_{ee}){sC_1} = 0
\]
\paragraph{Superposition:}
\begin{enumerate}
  \item Due to \(V_{in}\) alone (\(V_{ee} = 0V\)):
    \begin{align*}
      (V_{out} - V_{in}){G_1} + (V_{out} - 0){sC_1} &= 0 \\
      V_{out}G_1 - V_{in}G_1 + V_{out}sC_1 &= 0 \\
      V_{out}(G_1 + sC_1) - V_{in}G_1 &= 0 \\
      V_{out}(G_1 + sC_1) &= V_{in}G_1 \\
      \boxed{H_1(s) = \frac{N_1(s)}{D_1(s)} = \frac{V_{out_1}}{V_{in}} = \frac{G_1}{G_1 + sC_1}}
    \end{align*}
  \item Due to \(V_{ee}\) alone (\(V_{in} = 0V\)):
    \begin{align*}
      (V_{out} - 0){G_1} + (V_{out} - V_{ee}){sC_1} &= 0 \\
      V_{out}G_1 + V_{out}sC_1 - V_{ee}sC_1 &= 0 \\
      V_{out}(G_1 + sC_1) - V_{ee}sC_1 &= 0 \\
      V_{out}(G_1 + sC_1) &= V_{ee}sC_1 \\
      \boxed{H_2(s) = \frac{N_2(s)}{D_2(s)} = \frac{V_{out_2}}{V_{ee}} = \frac{sC_1}{G_1 + sC_1}}
    \end{align*}
\end{enumerate}

\paragraph{Total Output:}
\begin{align*}
  V_{out} &= V_{out_1} + V_{out_2} \\
  V_{out} &= \frac{V_{in}G_1}{G_1 + sC_1} + \frac{V_{ee}sC_1}{G_1 + sC_1}
\end{align*}

\subsubsection{Analysis}
\begin{enumerate}
  \item \textbf{Poles (\(D(s) = 0\)):} \\
    For both \(H_1(s)\) and \(H_2(s)\):
      \begin{align*}
        D(s) \quad\rightarrow\quad G_1 + sC_1 = 0 \\
        s_p &= -\frac{G_1}{C_1} \\
      \end{align*}
    Negative real pole indactes a stable first-order system.
    Frequency corresponding to this pole:
    \begin{align*}
      s &= -2 \pi f_p \\
      -2 \pi f_p &= -\frac{G_1}{C_1} \\
      \boxed{f_p = \frac{G_1}{2 \pi C_1} = \frac{1}{2 \pi R_1 C_1}}
    \end{align*}
  \item \textbf{Zeroes (\(N(s) = 0\)):} \\
    For $H_1(s)$:
      \begin{align*}
        N_1(s) = G_1 = 0
      \end{align*}
    No zeroes for \(H_1(s)\). \\
    For $H_2(s)$:
      \begin{align*}
        N_2(s) = sC_1 = 0 \\
        s_z &= 0 \quad \text{[let \(s = s_z\)]}\\
        f_z &= \frac{|s_z|}{2\pi} = 0 Hz
      \end{align*}
    Zero at origin for \(H_2(s)\).
  \item \textbf{DC Gain (\(H(0)\)):} \\
    For \(H_1(s)\):
      \begin{align*}
        H_1(0) &= \frac{G_1}{G_1 + 0 \cdot C_1} = 1 \\
        \boxed{A_{DC_1} = 1 = 0 dB}
      \end{align*}
    Unity DC gain for \(H_1(s)\). \\
    For \(H_2(s)\):
      \begin{align*}
        H_2(0) &= \frac{0 \cdot C_1}{G_1 + 0 \cdot C_1} = 0 \\
        \boxed{A_{DC_2} = 0 = -\infty dB}
      \end{align*}
    Zero DC gain for \(H_2(s)\).
  \item \textbf{Gain-Bandwidth Product (\(GBW\)):}
    We use the values from \(H_1(s)\) as it is the
    low-pass characteristic:
    \begin{align*}
      \text{GBW} = |A_{DC_1} \times f_p| = \frac{1}{2\pi R_1C_1}
    \end{align*}
  \item \textbf{Impedances:}
    Substituting \(G_1 = \frac{1}{R_1}\) and \(s = j \omega\):
    \paragraph{At Input:}
      \begin{align*}
        Z_{in} = R_1 + \frac{1}{sC_1} \\
        \text{Substituting } s = j \omega &: \\
        \boxed{Z_{in} = R_1 - j \frac{1}{\omega C_1}}
      \end{align*}
    \paragraph{At Output:}
      \begin{align*}
        Z_{out} &= R1 \parallel \frac{1}{sC_1} \\
        Z_{out} &= \frac{R1 \cdot \frac{1}{sC_1}}{R1 + \frac{1}{sC_1}} \\
        Z_{out} &= \frac{\frac{R_1}{sC_1}}{\frac{sC_1 R_1 + 1}{sC_1}} \\
        Z_{out} &= \frac{R_1}{sC_1 R_1 + 1} \\
        \text{Substituting } s = j \omega &: \\
        Z_{out} &= \frac{R_1}{1 + j \omega R_1 C_1} \cdot \frac{1 - j \omega R_1 C_1}{1 - j \omega R_1 C_1} \\
        Z_{out} &= \frac{R_1 - j \omega R_1^2 C_1}{1 + (\omega R_1 C_1)^2} \\
        \boxed{Z_{out} = \frac{R_1}{1 + (\omega R_1 C_1)^2} - j \frac{\omega R_1^2 C_1}{1 + (\omega R_1 C_1)^2}}
      \end{align*}
\end{enumerate}

\subsubsection{Question}
Design a RC Low Pass Filter with a cutoff frequency of
\(f_p = 45kHz\) using a capacitor value of \(C_1 = 1\mu F\).
Also compute the input and output impedances at
\(f = 10kHz\).
\paragraph{Solution:}
\begin{enumerate}
  \item Calculating \(R_1\):
    \begin{align*}
      f_p &= \frac{1}{2 \pi R_1 C_1} \\
      R_1 &= \frac{1}{2 \pi f_p C_1} \\
      &= \frac{1}{2 \pi (45kHz)(1\mu F)} \\
      R_1 &\approx 3.54 \Omega
    \end{align*}
  \item Calculating \(Z_{in}\) and \(Z_{out}\) at \(f = 10kHz\):
    \begin{align*}
      \omega &= 2 \pi f = 2 \pi (10kHz) = 62,832 \text{ rad/s} \\
      Z_{in} &= R_1 - j \frac{1}{\omega C_1} \\
             &= 3.54 \Omega - j \frac{1}{(62,832)(1\mu F)} \\
             &\approx 3.54 \Omega - j 15.9 \Omega \\
      Z_{out} &= \frac{R_1}{1 + (\omega R_1 C_1)^2} - j \frac{\omega R_1^2 C_1}{1 + (\omega R_1 C_1)^2} \\
              &= \frac{3.54 \Omega}{1 + (62,832 \cdot 3.54 \Omega \cdot 1\mu F)^2} - j \frac{(62,832)(3.54 \Omega)^2 (1\mu F)}{1 + (62,832 \cdot 3.54 \Omega \cdot 1\mu F)^2} \\
              &= \frac{3.54 \Omega}{1 + (0.222)^2} - j \frac{(62,832)(12.52 \Omega^2)(1\mu F)}{1 + (0.222)^2} \\
              &\approx \frac{3.54 \Omega}{1.049} - j \frac{0.786 \Omega}{1.049} \\
              &\approx 3.37 \Omega - j 0.75 \Omega
    \end{align*}
\end{enumerate}
